% grlite_report_v1.tex
% GRlite Technical Report, v1 -- "the receipts".
%
% A consolidated, three-part technical report that supersedes the incremental
% gr_sandbox_vNN design log:
%   Part I   -- the physics (derived in 3D, specialized to 2D as we go);
%   Part II  -- the PIC numerical scheme built from the 2D equations;
%   Part III -- the C library that embodies the scheme (API + correspondence).
%
% Built incrementally; this file currently contains the front matter and Part I.
\documentclass[11pt]{article}

\usepackage[T1]{fontenc}
\usepackage[utf8]{inputenc}
\usepackage{lmodern}
\usepackage[margin=1in]{geometry}
\usepackage{amsmath,amssymb}
\usepackage{xcolor}
\usepackage{array}
\usepackage{booktabs}
\usepackage{mdframed}
\usepackage{enumitem}
\usepackage[hidelinks]{hyperref}

% --- Notation macros (kept compatible with the gr_sandbox_vNN lineage) -------
\newcommand{\Phig}{\Phi_g}                 % gravitoelectric scalar potential
\newcommand{\Ag}{\mathbf{A}_g}             % gravitomagnetic vector potential
\newcommand{\vv}[1]{\mathbf{#1}}           % 3- or 2-vector
\newcommand{\dd}{\mathrm{d}}               % differential
\newcommand{\dx}{\Delta x}                 % grid spacing
\newcommand{\dt}{\Delta t}                 % time step
\newcommand{\order}[1]{\mathcal{O}(#1)}    % order symbol
\newcommand{\Geff}{G_{\mathrm{eff}}}
\newcommand{\keff}{k_e}
\newcommand{\ceff}{c}                       % c denotes the effective wave speed c_eff
\newcommand{\clocal}{c_{\mathrm{local}}}
\newcommand{\Bg}{\vv{B}_g}

% Left-aligned fixed-width table column.
\newcolumntype{L}[1]{>{\raggedright\arraybackslash}p{#1}}

% Function-name column: like L, but lets long monospace names break at their
% underscores so they wrap inside the cell instead of running into the next
% column or past the right margin (used for the API reference tables, S19).
\newcommand{\fnbreaks}{\renewcommand{\_}{\textunderscore\discretionary{}{}{}}}
\newcolumntype{F}[1]{>{\raggedright\arraybackslash\fnbreaks}p{#1}}

% Let inline monospace function names break at their underscores too, so long
% names in body text wrap instead of overrunning the right margin.  We patch the
% inner worker of the (robust) \texttt so it stays robust for moving arguments
% (ToC / PDF bookmarks); the added breaks are discretionary, so ordinary lines
% are unaffected.
\expandafter\let\expandafter\origtexttt\csname texttt \endcsname
\expandafter\def\csname texttt \endcsname#1{{\fnbreaks\origtexttt{#1}}}

% Code listings one size down, so the wider struct lines fit the text width.
\makeatletter
\renewcommand{\verbatim@font}{\normalfont\ttfamily\small}
\makeatother

% Boxed display ("the result you implement").
\newmdenv[linewidth=0.6pt,roundcorner=2pt,backgroundcolor=black!3,
          innertopmargin=4pt,innerbottommargin=4pt]{eqbox}

% Algorithm step list ("Step N" referenced in prose).
\newlist{algsteps}{enumerate}{1}
\setlist[algsteps]{label=\textbf{Step \arabic*.},leftmargin=2.6em,itemsep=4pt}

% 2D-specialization callout.
\newmdenv[linewidth=0.6pt,roundcorner=2pt,linecolor=black!40,
          backgroundcolor=blue!4,innertopmargin=4pt,innerbottommargin=4pt,
          skipabove=6pt,skipbelow=6pt]{twoD}

\title{\textbf{GRlite}\\[2pt]
  \large A Real-Time 2D Linearized-Gravity + Electromagnetism\\
  Particle-in-Cell Sandbox\\[6pt]
  \normalsize Technical Report v1}
\author{D.\ Marks\\\small\texttt{github.com/profdc9/GRlite}}
\date{\today}

\begin{document}
\maketitle

\begin{abstract}
GRlite is an interactive sandbox that integrates \emph{linearized} general
relativity coupled to electromagnetism on a 2D grid, in real time, in a web
browser.  It is a particle-in-cell (PIC) field solver: massive and charged
macroparticles source six gravito-electromagnetic potentials on a Yee lattice,
the potentials are advanced by an explicit wave-equation leapfrog, and the
particles are pushed by a relativistic force law gathered from those potentials.
This report documents the simulator at the level of detail of a technical
report, in three parts.  \textbf{Part~I} derives the physical model: linearized
GR in the gravito-electromagnetic (GEM) form, its coupling to Maxwell
electromagnetism, the relativistic single-particle dynamics (including
post-Newtonian corrections, spin, and proper time), and the
background/perturbation split.  Each result is derived in 3D and then
specialized to the $2{+}1$ dimensions the simulator actually runs in, so the
correspondence between the full physics and the reduced model is explicit.
\textbf{Part~II} (separate) builds the PIC scheme from the 2D equations: Yee
sublattices and the leapfrog field update; the CIC/TSC/bump shape functions;
charge-conserving Esirkepov current deposition; the Lewis--Birdsall
momentum-conserving force gather; the relativistic Boris pusher; initial-field
generation; the absorbing boundary; and the visualization mathematics.
\textbf{Part~III} (separate) documents the C library that implements the scheme
and maps every public function to the step of the PIC algorithm it realizes.
\end{abstract}

\tableofcontents

%====================================================================
\section*{Conventions, units, and notation}
\addcontentsline{toc}{section}{Conventions, units, and notation}
%====================================================================

\paragraph{Signature and indices.}
The metric signature is $(-,+,+,+)$.  Greek indices $\mu,\nu,\dots$ run over
$0,1,2,3$ (spacetime); Latin indices $i,j,k$ over the spatial dimensions.  Bold
symbols ($\vv{x}$, $\vv{v}$, $\Ag$) are spatial vectors --- 3-vectors in the
derivations of this Part, 2-vectors after the 2D reduction.

\paragraph{Effective constants.}
The simulator does not use SI values for the fundamental constants; it exposes
\emph{tunable} effective constants so a student can dial gravity, electrostatics,
and the wave speed independently and explore limits ($\ceff\to\infty$ recovers
the Newtonian limit; $\Geff\to 0$ switches off gravity; etc.).  Throughout this
report the symbols $G$, $c$, and the Coulomb constant $\keff = 1/(4\pi
\varepsilon_0)$ denote those effective quantities $\Geff$, $c_{\mathrm{eff}}$,
$\keff$.  We write $G$ and $c$ for readability and flag the substitution where
it matters numerically (e.g.\ the Schwarzschild precession rate
$\propto \Geff M / c_{\mathrm{eff}}^2$).

\paragraph{The two object classes.}
GRlite distinguishes two kinds of gravitating/charged object:
\emph{macroparticles}, which move, deposit fields on the grid, and feel forces;
and \emph{background bodies}, which are fixed analytic sources (a precomputed
field that never moves and feels no back-reaction).  Part~I treats their common
physics; the split itself is the subject of \S\ref{sec:bgsplit}.

\paragraph{The six potentials.}
The dynamical state of the field is six scalar potentials,
\begin{equation}
  \underbrace{\Phig,\; A_{g,x},\; A_{g,y}}_{\text{gravity (GEM)}}\,,\qquad
  \underbrace{\phi,\; A_x,\; A_y}_{\text{electromagnetism}}\,,
\end{equation}
each advanced by a sourced wave equation.  Derived fields ($\vv{g}$, $\Bg$,
$\vv{E}$, $\vv{B}$) are never stored; they are reconstructed on demand from
finite differences of these potentials.  This ``potentials are the state''
choice drives much of both the physics formulation (Part~I) and the numerics
(Part~II).

\part{Physical Foundations}

%====================================================================
\section{Linearized General Relativity and Gravito-electromagnetism}
\label{sec:lingr}
%====================================================================

\subsection{Goal and regime}

A full numerical-relativity evolution (e.g.\ the ADM $3{+}1$ split with a Cauchy
boundary) is far too expensive for an interactive desktop/browser tool.  GRlite
instead works in \emph{linearized} GR, where the metric is a small perturbation
of flat spacetime,
\begin{equation}
  g_{\mu\nu} = \eta_{\mu\nu} + h_{\mu\nu}, \qquad |h_{\mu\nu}| \ll 1 .
\end{equation}
In the Lorenz gauge $\partial_\mu \bar h^{\mu\nu}=0$, the trace-reversed
perturbation $\bar h^{\mu\nu} \equiv h^{\mu\nu} - \tfrac12 \eta^{\mu\nu} h$
obeys a sourced wave equation,
\begin{equation}
\label{eq:hbar_wave}
  \Box\, \bar h^{\mu\nu} = -\frac{16\pi G}{c^4}\, T^{\mu\nu}_{\mathrm{total}},
  \qquad \Box \equiv -\frac{1}{c^2}\partial_t^2 + \nabla^2 .
\end{equation}
This is structurally identical to the wave equation for the electromagnetic
4-potential, which is exactly why FDTD methods from computational
electromagnetics carry over directly --- and why the whole simulator can treat
gravity and electromagnetism with one field engine.

\subsection{The GEM metric}
\label{sec:gem_metric}

Choosing the gravito-electromagnetic (GEM) potentials $\Phig$ (gravitoelectric
scalar) and $\Ag$ (gravitomagnetic vector) to carry the relevant components of
$\bar h^{\mu\nu}$, the linearized metric is
\begin{equation}
\label{eq:gem_metric}
  h_{00} = -\frac{2\Phig}{c^2}, \qquad
  h_{0i} = \frac{2 A_{g,i}}{c}, \qquad
  h_{ij} = -\frac{2\Phig}{c^2}\,\delta_{ij},
\end{equation}
with line element
\begin{equation}
\label{eq:gem_line}
  \dd s^2 = -\Bigl(1+\tfrac{2\Phig}{c^2}\Bigr)c^2\,\dd t^2
            + \frac{4}{c}\,\Ag\cdot\dd\vv{x}\,\dd t
            + \Bigl(1-\tfrac{2\Phig}{c^2}\Bigr)\bigl(\dd x^2+\dd y^2+\dd z^2\bigr).
\end{equation}
The trace-reversed components are $\bar h_{00} = -4\Phig/c^2$,
$\bar h_{0i} = 4 A_{g,i}/c$, $\bar h_{ij}=0$, so
Eq.~\eqref{eq:hbar_wave} decouples cleanly into
$\Box\Phig \propto \rho$ and $\Box\Ag \propto \vv{J}$ with no spatial-stress
source at this order.

\paragraph{What each metric component does.}
The three distinct pieces of Eq.~\eqref{eq:gem_metric} are responsible for
disjoint classes of effect, and the simulator can switch them independently:
\begin{itemize}
  \item $h_{00}=-2\Phig/c^2$ (time--time): the Newtonian potential. Gives
  gravitational attraction $-\nabla\Phig$, gravitational time dilation
  ($+\Phig/c^2$ in $\dd\tau/\dd t$), the Shapiro delay, and \emph{half} of
  light bending.
  \item $h_{ij}=-2\Phig/c^2\,\delta_{ij}$ (spatial): no Newtonian analogue and
  invisible to slow particles, but supplies the \emph{second half} of light
  bending (recovering the GR $4GM/c^2 b$), the post-Newtonian force corrections
  (\S\ref{sec:relcorr}), and hence perihelion precession.
  \item $h_{0i}=2A_{g,i}/c$ (time--space): nonzero only for moving/rotating
  mass currents. Produces the gravitomagnetic force, Lense--Thirring gyroscope
  precession $\boldsymbol\Omega_{\mathrm{LT}}=\nabla\times\Ag$, the
  gravitomagnetic clock effect (\S\ref{sec:proptime}), and frame dragging.
\end{itemize}

\subsection{GEM field equations and the Lorentz force}

In the slow-motion weak-field limit the linearized equations take Maxwell form
for a gravitoelectric field $\vv{g}$ and gravitomagnetic field $\Bg$:
\begin{align}
  \nabla\cdot\vv{g}   &= -4\pi G\rho, &
  \nabla\times\vv{g}  &= -\partial_t\Bg, \\
  \nabla\cdot\Bg      &= 0, &
  \nabla\times\Bg     &= -\frac{4\pi G}{c^2}\vv{J} + \frac{1}{c^2}\partial_t\vv{g},
\end{align}
with $\vv{g}=-\nabla\Phig-\partial_t\Ag$ and $\Bg=\nabla\times\Ag$.  A test
particle of rest mass $m$ and charge $q$ feels the combined GEM+EM Lorentz
force
\begin{equation}
\label{eq:lorentz_fields}
  \vv{F} = m\bigl(\vv{g} + 4\,\vv{v}\times\Bg\bigr)
         + q\bigl(\vv{E} + \vv{v}\times\vv{B}\bigr).
\end{equation}
The factor of $4$ on the gravitomagnetic term is the fingerprint of the spin-2
graviton (versus spin-1 photon) and is responsible for the Lense--Thirring
rate; getting it right also fixes the metric coefficient $h_{0i}=2A_{g,i}/c$ and
the proper-time cross term (\S\ref{sec:proptime}).

\subsection{Potential-form force: what the code actually evaluates}
\label{sec:force_potential}

Because the simulator stores \emph{potentials}, not fields, it is essential to
rewrite Eq.~\eqref{eq:lorentz_fields} so that no grid-wide curl is required.
Using $\vv{v}\times(\nabla\times\vv{A}) = \nabla(\vv{v}\cdot\vv{A}) -
(\vv{v}\cdot\nabla)\vv{A}$,
\begin{align}
  \vv{F}_{\mathrm{GEM}} &= m\bigl(-\nabla\Phig - \partial_t\Ag
     + 4\nabla(\vv{v}\cdot\Ag) - 4(\vv{v}\cdot\nabla)\Ag\bigr),
     \label{eq:force_gem_pot}\\
  \vv{F}_{\mathrm{EM}}  &= q\bigl(-\nabla\phi - \partial_t\vv{A}
     + \nabla(\vv{v}\cdot\vv{A}) - (\vv{v}\cdot\nabla)\vv{A}\bigr).
     \label{eq:force_em_pot}
\end{align}
Every term is a spatial or temporal derivative of a stored potential evaluated
at the particle. The force is gauge-invariant: under
$\phi\to\phi-\partial_t\Lambda,\ \vv{A}\to\vv{A}+\nabla\Lambda$ the per-term
gauge dependence cancels. These two expressions (with the relativistic
corrections of \S\ref{sec:relcorr} added to the gravity piece) are the force law
the pusher of Part~II implements.

\begin{twoD}
\textbf{2D specialization.} In $2{+}1$ dimensions $\vv{g}$ and $\vv{E}$ are
2-vectors; $\Bg$ and $\vv{B}$ collapse to the single scalar curl components
$B_{g,z}=\partial_x A_{g,y}-\partial_y A_{g,x}$ and
$B_z=\partial_x A_y-\partial_y A_x$. Cross products become scalar multiplies,
$(\vv{v}\times\Bg)=(v_y B_{g,z},\,-v_x B_{g,z})$, so the field-form force
\eqref{eq:lorentz_fields} reads
$F_x = m(g_x+4v_yB_{g,z})+q(E_x+v_yB_z)$,
$F_y = m(g_y-4v_xB_{g,z})+q(E_y-v_xB_z)$.
The potential-form expressions \eqref{eq:force_gem_pot}--\eqref{eq:force_em_pot}
are unchanged in form with $\nabla=(\partial_x,\partial_y)$.
\end{twoD}

\subsection{Particle properties: mass and charge are independent}

Each macroparticle carries a rest mass $m$ and an electric charge $q$ that
couple to disjoint fields:
\begin{itemize}
  \item rest mass $m$ sources gravity, $\rho_{\mathrm{matter}}=\sum_i\gamma
  m_i\,\delta(\vv{x}-\vv{x}_i)$ and $\vv{J}_{\mathrm{matter}}=\sum_i\gamma
  m_i\vv{v}_i\,\delta(\vv{x}-\vv{x}_i)$, and sets the inertia via
  $\vv{p}=\gamma m\vv{v}$;
  \item charge $q$ sources EM, $\rho_q=\sum_i q_i\,\delta(\vv{x}-\vv{x}_i)$ and
  $\vv{J}_q=\sum_i q_i\vv{v}_i\,\delta(\vv{x}-\vv{x}_i)$, and sets the response
  to $q(\vv{E}+\vv{v}\times\vv{B})$.
\end{itemize}
A particle with both $m\neq0$ and $q\neq0$ contributes to all four potential
families and feels all four force families.

%====================================================================
\section{Electromagnetism in the Linearized Background}
\label{sec:em}
%====================================================================

\subsection{Potentials and wave equations}

Maxwell's equations in the Lorenz gauge give the same wave-equation structure as
gravity,
\begin{equation}
  \Box\phi = -\rho_q/\varepsilon_0, \qquad \Box\vv{A} = -\mu_0\vv{J}_q,
\end{equation}
with $\vv{E}=-\nabla\phi-\partial_t\vv{A}$ and $\vv{B}=\nabla\times\vv{A}$.  In
the GRlite engine these are advanced by the \emph{same} leapfrog as the GEM
potentials; the only difference is the source and the local wave speed below.

\subsection{Shapiro delay and gravitational lensing via \texorpdfstring{$\clocal$}{c\_local}}
\label{sec:shapiro}

A photon follows a null geodesic, $\dd s^2=0$.  Substituting the line element
\eqref{eq:gem_line} (with both the time and the spatial coefficient) and solving
for the coordinate light speed gives
\begin{equation}
\label{eq:clocal}
  \frac{\dd x}{\dd t} \approx c\Bigl(1 + \frac{2\Phig}{c^2}\Bigr) \equiv \clocal .
\end{equation}
The factor of $2$ comes from both metric coefficients contributing to the null
condition; it is the same factor that doubles light bending over the Newtonian
value.  The Shapiro time delay is $\Delta t = -\tfrac{2}{c^3}\int_{\rm
path}\Phig\,\dd\ell$.  GRlite realizes both Shapiro delay and lensing by
\emph{stepping the EM wave equation with $\clocal(\vv{x})$ in place of $c$};
wavefronts slow and bend in the potential well exactly as light should.
Crucially $\Phig$ here is the \emph{total} (background + perturbation) potential,
so a moving deposited mass can lens an EM wave (a dynamic lens).

\subsection{Electromagnetic stress-energy as a gravitational source}
\label{sec:energy}

Electromagnetic field energy gravitates. Writing $\vv{P}\equiv\nabla\phi +
\partial_t\vv{A}=-\vv{E}$ and $\vv{C}\equiv\nabla\times\vv{A}=\vv{B}$, the EM
contributions to the GEM source ($T^{00}$ energy density and $T^{0i}$ Poynting
flux, divided by $c^2$) are
\begin{align}
  \rho_{\mathrm{EM}} &= \frac{\varepsilon_0}{2c^2}\bigl(|\vv{P}^{\rm pert}|^2 +
     c^2|\vv{C}^{\rm pert}|^2\bigr), \label{eq:rho_em}\\
  \vv{J}_{\mathrm{EM}} &= -\frac{1}{\mu_0 c^2}\,\vv{P}^{\rm pert}\times\vv{C}^{\rm pert}.
     \label{eq:J_em}
\end{align}
Only the \emph{perturbation} potentials enter, to avoid double-counting the
background metric.  The Maxwell spatial stress $T^{ij}$ is dropped (it would
source $h^{ij}$, which the GEM reduction does not evolve); this is the principal
approximation of the ``effective-source'' method and is valid when EM fields are
not so strong or directional that $T^{ij}$ matters
($|\vv{S}|/c \ll u_{\mathrm{EM}}$).

\begin{twoD}
\textbf{2D specialization.} The curl becomes the scalar
$F_{12}^{\rm pert}=\partial_x A_y^{\rm pert}-\partial_y A_x^{\rm pert}$, and the
cross product in \eqref{eq:J_em} becomes a $90^\circ$ rotation of the 2-vector
$\vv{P}^{\rm pert}$ scaled by $F_{12}^{\rm pert}$:
\begin{equation}
  \rho_{\mathrm{EM}} = \frac{\varepsilon_0}{2c^2}\bigl(|\vv{P}^{\rm pert}|^2 +
     c^2 (F_{12}^{\rm pert})^2\bigr),\quad
  \vv{J}_{\mathrm{EM}} = -\frac{F_{12}^{\rm pert}}{\mu_0 c^2}
     \begin{pmatrix} P_y^{\rm pert}\\ -P_x^{\rm pert}\end{pmatrix}.
\end{equation}
\end{twoD}

%====================================================================
\section{The Unified Six-Field System}
\label{sec:sixfield}
%====================================================================

Collecting \S\S\ref{sec:lingr}--\ref{sec:em}, the dynamical field state and its
sources are summarized in Table~\ref{tab:sixfield}.  All six potentials obey the
\emph{same} wave operator; gravity and EM differ only in coupling sign/strength
and (for EM) the use of $\clocal$.  This uniformity is what lets a single
leapfrog kernel (Part~II) advance the entire system.

\begin{table}[h]
\centering
\begin{tabular}{@{}L{2.4cm}L{4.6cm}L{2.6cm}L{3.2cm}@{}}
\toprule
\textbf{Potential} & \textbf{Wave equation} & \textbf{Source} & \textbf{Derived field} \\
\midrule
$\Phig$    & $\Box\Phig = -4\pi G\rho_{\rm eff}$              & $\rho_{\rm eff}$            & $\vv{g}=-\nabla\Phig-\partial_t\Ag$ \\
$A_{g,x},A_{g,y}$ & $\Box\Ag = -\tfrac{4\pi G}{c^2}\vv{J}_{\rm eff}$ & $\vv{J}_{\rm eff}$ & $B_{g,z}=(\nabla\times\Ag)_z$ \\
$\phi$     & $\Box_{\clocal}\phi = -\rho_q/\varepsilon_0$     & $\rho_q$                   & $\vv{E}=-\nabla\phi-\partial_t\vv{A}$ \\
$A_x,A_y$  & $\Box_{\clocal}\vv{A} = -\mu_0\vv{J}_q$          & $\vv{J}_q$                 & $B_z=(\nabla\times\vv{A})_z$ \\
\bottomrule
\end{tabular}
\caption{The six-potential system. $\rho_{\rm eff}=\rho_{\rm matter}+\rho_{\rm
EM}$ and $\vv{J}_{\rm eff}=\vv{J}_{\rm matter}+\vv{J}_{\rm EM}$ fold the EM
stress-energy \eqref{eq:rho_em}--\eqref{eq:J_em} back into the gravitational
source. In 2D the two vector potentials per family are genuine 2-vectors and the
curls are scalars.}
\label{tab:sixfield}
\end{table}

%====================================================================
\section{Relativistic Particle Dynamics}
\label{sec:dynamics}
%====================================================================

\subsection{Relativistic momentum update}

Particles are integrated through the relativistic 3-momentum $\vv{p}=\gamma
m\vv{v}$, $\gamma = (1-v^2/c^2)^{-1/2}$:
\begin{equation}
\label{eq:pupdate}
  \vv{p}\leftarrow \vv{p}+\vv{F}_{\rm total}\,\dd t, \qquad
  \vv{v} = \frac{\vv{p}}{\sqrt{m^2+|\vv{p}|^2/c^2}}, \qquad
  \vv{x}\leftarrow\vv{x}+\vv{v}\,\dd t .
\end{equation}
Pushing $\vv{p}$ (not $\vv{v}$) automatically captures the anisotropy of
relativistic inertia (longitudinal mass $\gamma^3 m$ vs.\ transverse $\gamma
m$).  The EM force \eqref{eq:force_em_pot} is already exactly relativistic; the
gravitational force needs the corrections below.

\subsection{Post-Newtonian (EIH) corrections}
\label{sec:relcorr}

GEM as written is a slow-motion expansion: at relativistic speed the spatial
metric $h^{ij}$ also deflects particles (this is the missing ``second half'' of
light bending).  The correct $\order{(v/c)^2}$ equation of motion is the
test-particle limit of the Einstein--Infeld--Hoffmann (EIH) 1PN
equation~\cite{Will1993,AliHaimoud2019}.  Writing $\vv{g}=-\nabla\Phig-\partial_t\Ag$
and organizing by order in $v/c$,
\begin{equation}
\label{eq:geodesic_expansion}
  \vv{a} = \underbrace{\vv{g}}_{\order{1}}
        + \underbrace{4\,\vv{v}\times\Bg}_{\order{v/c}\ \text{(Lense--Thirring)}}
        + \underbrace{\frac{v^2}{c^2}\vv{g} + \frac{4\Phig}{c^2}\vv{g}
           - \frac{4(\vv{v}\cdot\vv{g})\vv{v}}{c^2}}_{\order{v^2/c^2}\ \text{(EIH 1PN)}}
        + \order{v^3/c^3}.
\end{equation}
The three 1PN terms are physically distinct: $(v^2/c^2)\vv{g}$ strengthens
deflection of fast particles (and, with the $\clocal$ spatial-metric piece,
recovers the factor-of-2 photon deflection as $v\to c$); $(4\Phig/c^2)\vv{g}$ is
a radial ``Shapiro'' enhancement contributing to binding energy and precession;
and $-4(\vv{v}\cdot\vv{g})\vv{v}/c^2$ is a velocity-coupling term that does no
net work but breaks Keplerian closure.  Integrated over a bound orbit they
reproduce the Schwarzschild perihelion advance
\begin{equation}
\label{eq:precession}
  \Delta\varphi_{\rm prec} = \frac{6\pi\,\Geff M}{c_{\rm eff}^2\,a(1-e^2)} .
\end{equation}
The simulator exposes these as force \emph{tiers} (Newtonian; GEM; EIH 1PN),
selectable per scenario; Table~\ref{tab:tiers} lists them.  In implementation
the gravitomagnetic cross product is always replaced by its potential form
(\S\ref{sec:force_potential}), so $\Bg$ is never formed as a grid array; the
scalar $\vv{g}$ used in the 1PN terms is evaluated on the fly from the
potentials at the particle.

\begin{table}[h]
\centering
\begin{tabular}{@{}lL{8.2cm}l@{}}
\toprule
\textbf{Tier} & \textbf{Gravitational acceleration} & \textbf{Range} \\
\midrule
Newtonian & $\vv{g}$ & $v/c<0.01$ \\
GEM       & $\vv{g}+4\,\vv{v}\times\Bg$ (potential form) & $v/c<0.1$ \\
EIH 1PN   & GEM $+\ \tfrac{v^2}{c^2}\vv{g}+\tfrac{4\Phig}{c^2}\vv{g}-\tfrac{4(\vv{v}\cdot\vv{g})\vv{v}}{c^2}$ & $v/c\lesssim0.5$ \\
\bottomrule
\end{tabular}
\caption{Gravitational force tiers. EM is exactly relativistic in all tiers; the
momentum update \eqref{eq:pupdate} is always relativistic.}
\label{tab:tiers}
\end{table}

\subsection{Intrinsic spin}
\label{sec:spin}

A particle with intrinsic angular momentum obeys the
Mathisson--Papapetrou--Dixon equations; in the GEM/EM approximation these reduce
to a 3-vector spin $\vv{S}$ integrated alongside the momentum.  The spin obeys
the Tulczyjew--Dixon condition $S^{\mu\nu}p_\nu=0$ (note: $p^\mu$, the momentum,
not $mU^\mu$), whose 3+1 split fixes $S^0=\vv{S}\cdot\vv{p}/(\gamma m c)$ and is
maintained each step by the covariant projection
\begin{equation}
\label{eq:spin_constraint}
  \vv{S}\leftarrow \vv{S} - \frac{\vv{S}\cdot\vv{p}}{|\vv{p}|^2+m^2c^2}\,\vv{p}.
\end{equation}
Spin precesses under gravitomagnetic, EM, and Thomas terms,
\begin{equation}
\label{eq:precession_rates}
  \frac{\dd\vv{S}}{\dd t} = \vv{S}\times\Bigl(
     \underbrace{\nabla\times\Ag}_{\rm grav}
   + \underbrace{\tfrac{g_s q}{2m}\nabla\times\vv{A}}_{\rm EM}
   + \underbrace{-\tfrac{\gamma^2}{\gamma+1}\tfrac{\vv{v}\times\vv{a}}{c^2}}_{\rm Thomas}\Bigr),
\end{equation}
and, being a moving dipole, sources magnetization currents (magnetic moment
$\boldsymbol\mu=g_s(q/2m)\vv{S}$, gravitomagnetic moment $\boldsymbol\sigma =
2\vv{S}$). A spin gradient also exerts a Stern--Gerlach force
$\vv{F}_{\rm spin}=\nabla(\vv{S}\cdot\Bg)+\nabla(\boldsymbol\mu\cdot\vv{B})$.

\begin{twoD}
\textbf{2D specialization.} With motion in the plane and spin along $\hat z$,
$\vv{S}$ is a single scalar $s$ (an angle $\phi_{\rm spin}$), and
$\nabla\times\Ag,\ \nabla\times\vv{A}$ are the scalars $B_{g,z},\ B_z$.
Precession reduces to
$\dot\phi_{\rm spin}=B_{g,z}+\tfrac{g_s q}{2m}B_z+\Omega_{{\rm Thomas},z}$
with $\Omega_{{\rm Thomas},z}=-\tfrac{\gamma^2}{\gamma+1}(v_xa_y-v_ya_x)/c^2$,
the spin-gradient force becomes $s\,\nabla B_{g,z}+\mu\,\nabla B_z$, and the
Tulczyjew--Dixon constraint \eqref{eq:spin_constraint} is satisfied
automatically (no projection step needed).
\end{twoD}

\subsection{Proper time}
\label{sec:proptime}

Each particle integrates the proper time along its worldline from the metric
\eqref{eq:gem_line}, $\dd\tau^2=-\dd s^2/c^2$:
\begin{equation}
\label{eq:proper_time}
  \frac{\dd\tau}{\dd t} = \sqrt{\Bigl(1+\frac{2\Phig}{c^2}\Bigr)
     - \Bigl(1-\frac{2\Phig}{c^2}\Bigr)\frac{v^2}{c^2}
     - \frac{8\,\Ag\cdot\vv{v}}{c^2}}
  \;\approx\; 1 + \frac{\Phig}{c^2} - \frac{v^2}{2c^2} - \frac{4\,\Ag\cdot\vv{v}}{c^2}.
\end{equation}
The three first-order terms are the gravitational redshift ($\Phig/c^2$, slows
clocks in a well), the kinematic dilation ($-v^2/2c^2$), and the
\emph{gravitomagnetic clock effect} ($-4\Ag\cdot\vv{v}/c^2$).  The cross-term
coefficient is fixed to $8$ inside the root (equivalently $4$ at first order) by
consistency with the factor-of-4 force law / $h_{0i}=4A_{g,i}/c$ trace-reversed
component; this gives prograde clocks running \emph{slower} than retrograde at
equal radius, matching the Cohen--Mashhoon Kerr
result~\cite{CohenMashhoon1993,Mashhoon2003}. Around a closed orbit the
prograde$-$retrograde difference is the gravitomagnetic flux,
$\Delta\tau = -\tfrac{8}{c^2}\oint\Ag\cdot\dd\vv{x} =
-\tfrac{8}{c^2}\!\int\Bg\cdot\dd\vv{A}$.
The same interpolated scalar $\vv{v}\cdot\Ag$ already computed for the force is
reused here, so a proper-time clock costs one square root per particle per step.

%====================================================================
\section{Background Spacetime and the Perturbation Split}
\label{sec:bgsplit}
%====================================================================

\subsection{Decomposition and modes}

When fixed sources dominate the field, it is wasteful to evolve them on the grid.
GRlite splits every potential into a precomputed analytic \emph{background} and a
time-stepped \emph{perturbation},
\begin{equation}
  \Phig(\vv{x},t)=\Phig^{\rm bg}(\vv{x})+\Phig^{\rm pert}(\vv{x},t),
\end{equation}
and likewise for $\Ag,\phi,\vv{A}$.  The force on a particle uses the sum of
both contributions; only the perturbation is advanced by the wave solver, and it
is sourced only by the moving macroparticles.  Three modes follow: pure
background (test particles, zero field cost), background $+$ perturbation (a few
interacting particles in a dominant field), and fully dynamic (no background; all
fields evolved, e.g.\ an equal-mass binary).

\subsection{The compact body: a no-hair \texorpdfstring{$(M,Q,J)$}{(M,Q,J)} source}
\label{sec:compactbody}

The canonical background is a single fixed \emph{compact body} carrying the three
``no-hair'' charges of a stationary black hole --- mass $M$, electric charge
$Q$, angular momentum $J_z$ --- superposed in one object.  Its analytic
potentials are the softened point sources
\begin{align}
  \Phig^{\rm bg} &= -\frac{GM}{\sqrt{r^2+\varepsilon^2}}, &
  \phi^{\rm bg}  &= +\frac{\keff Q}{\sqrt{r^2+\varepsilon^2}}, \\
  \Ag^{\rm bg}   &= \frac{G J_z}{2c^2}\,\frac{\hat z\times\vv{r}}{(r^2+\varepsilon^2)^{3/2}}, &
  \vv{A}^{\rm bg} &= \frac{\keff\,\mu_z}{c^2}\,\frac{\hat z\times\vv{r}}{(r^2+\varepsilon^2)^{3/2}},
\end{align}
where $r=|\vv{x}-\vv{x}_0|$ and $\varepsilon$ is a softening length (a few
cells) that removes the singular core.  The special cases reproduce the
weak-field analogues of the classical metrics: $(M,0,0)$ Schwarzschild,
$(M,Q,0)$ Reissner--Nordstr\"om, $(M,0,J)$ Kerr, $(M,Q,J)$ Kerr--Newman.

\paragraph{The magnetic dipole is fixed by the no-hair relation.}
A spinning charge has an EM magnetic dipole moment set by the gyromagnetic
relation
\begin{equation}
\label{eq:nohair_mu}
  \mu_z = g_s\,\frac{Q}{2M}\,J_z
        = \frac{g_s\,Q\,J_z\,\Geff}{2\,GM},
\end{equation}
with $g_s=2$ reproducing the Kerr--Newman value $\mu=QJ/M$.  Equation
\eqref{eq:nohair_mu} is undefined for $M=0$, so a massless charged rotor carries
\emph{no} EM dipole field --- the correct degenerate limit, since such an object
is not a Kerr--Newman body.  (The gravitomagnetic dipole $\Ag^{\rm bg}$, by
contrast, is sourced by $J_z$ alone and is mass-independent.)

\paragraph{Superposition.}
Multiple compact bodies superpose linearly: the background potential is the sum
of each body's analytic field. This is exact in the linearized theory (the field
equations are linear), and is the basis for multi-body scenes --- the
two-fixed-center problem, multi-lens gravitational lensing, and mixed
$M/Q/J$ fields. The simulator supports up to a fixed number of such bodies
(Part~III).

\subsection{The 2D Green's function choice}
\label{sec:greens2d}

In strict 2D the static Green's function of $\nabla^2$ is logarithmic, so a 2D
point mass gives $\Phig\sim\ln r$ and $|\vv{g}|\sim 1/r$ (flat rotation curves)
rather than the 3D $\Phig\sim 1/r$, $|\vv{g}|\sim1/r^2$:
\begin{equation}
  \text{3D: } \Phig\sim 1/r,\ |\vv{g}|\sim1/r^2; \qquad
  \text{2D: } \Phig\sim\ln r,\ |\vv{g}|\sim1/r .
\end{equation}
GRlite makes a deliberate pedagogical choice here: the \emph{deposited}
(perturbation) field is the honest 2D solver output (a $\ln r$ field, since it
solves the 2D Poisson/wave equation), while the analytic \emph{background} body
prescribes the 3D-flavored softened $1/r$ form of \S\ref{sec:compactbody}, so
that test-particle orbits look Keplerian and the precession test
\eqref{eq:precession} is meaningful. The FDTD perturbation solver is identical
in either case; only the analytic background differs. This dual nature (true-2D
deposited fields vs.\ 3D-flavored fixed sources) is intrinsic to the model and
is flagged again where it affects the numerics (Part~II) and the visualizations.

%====================================================================
\section{What Survives in 2D: Summary}
\label{sec:2dsummary}
%====================================================================

Table~\ref{tab:2d} consolidates the 2D reduction used throughout Part~I and
assumed by Parts~II--III. The headline is that essentially all of the
physically interesting GR/EM effects survive the reduction; what is lost is the
transverse tensor (gravitational-wave $+/\times$) polarizations, which require
two transverse spatial directions absent in $2{+}1$.

\begin{table}[h]
\centering
\begin{tabular}{@{}L{4.6cm}cL{6.4cm}@{}}
\toprule
\textbf{Quantity / effect} & \textbf{2D?} & \textbf{Reduced form / note} \\
\midrule
$\vv{g},\vv{E}$               & yes & 2-vectors $(x,y)$ \\
$\Bg,\vv{B}$                  & yes & scalars $B_{g,z},B_z$ (the curl) \\
$\vv{v}\times\Bg$             & yes & $(v_yB_{g,z},-v_xB_{g,z})$ \\
spin $\vv{S}$                 & yes & scalar angle $\phi_{\rm spin}$; TD constraint automatic \\
Newtonian gravity            & yes & $\ln r$ deposited; $1/r$ available as background \\
frame dragging               & yes & $B_{g,z}$ from spin/moving mass \\
perihelion precession        & yes & $\order{v^2/c^2}$ EIH terms \\
Shapiro delay / lensing      & yes & via $\clocal$ \\
proper-time dilation         & yes & per-particle clock, Eq.~\eqref{eq:proper_time} \\
gravitomagnetic clock effect & yes & $-8\Ag\cdot\vv{v}/c^2$ term \\
EM--gravity coupling         & yes & $\rho_{\rm EM},\vv{J}_{\rm EM}$ \\
GW scalar mode               & yes & propagates at $c$ \\
GW tensor ($+,\times$) modes & \textbf{no} & require $h^{ij}$; absent in 2D GEM \\
\bottomrule
\end{tabular}
\caption{The $2{+}1$ reduction at a glance.}
\label{tab:2d}
\end{table}

\bigskip
\noindent\textit{End of Part~I.}  Part~II builds the particle-in-cell numerical
scheme directly on the 2D equations collected in
Table~\ref{tab:sixfield} and \S\ref{sec:2dsummary}; Part~III documents the C
library that implements it.

\part{The Particle-in-Cell Scheme}

%====================================================================
\section{Overview: the per-step pipeline}
\label{sec:pipeline}
%====================================================================

GRlite is a particle-in-cell (PIC) code: the continuous sources of
Part~I are represented by a finite set of macroparticles, deposited onto grid
arrays each step; the six potentials are advanced by an explicit wave-equation
leapfrog; and the particles are pushed by the force gathered back from those
potentials. The dynamical state is

\begin{itemize}\setlength\itemsep{2pt}
  \item six \textbf{potential} arrays $\{\Phig,A_{g,x},A_{g,y},\phi,A_x,A_y\}$,
  each carrying three time levels (prev $=\,$\textsuperscript{$n-1$},
  curr $=\,$\textsuperscript{$n$}, next $=\,$\textsuperscript{$n+1$}) on its Yee
  sublattice (\S\ref{sec:lattice});
  \item six \textbf{source} arrays
  $\{\rho_{\rm matter},J_{m,x},J_{m,y},\rho_q,J_{q,x},J_{q,y}\}$ rebuilt each
  step from the particles;
  \item a \textbf{particle} array, each entry storing
  $(x,y,p_x,p_y,m,q,\tau,\phi_{\rm spin},g_s,\dots)$;
  \item optional fixed \textbf{background} fields (\S\ref{sec:bgimpl}).
\end{itemize}

One time step (the routine \texttt{gr\_sim\_step}, Part~III) executes:

\begin{algsteps}
  \item \textbf{Deposit sources.} Clear the source arrays; for each particle
  deposit $\rho^{\,n}$ at $\vv{x}^n$ with the chosen shape (\S\ref{sec:shapes}),
  the current $\vv{J}^{\,n-1/2}$ over the segment $\vv{x}^{n-1}\!\to\vv{x}^n$ by
  Esirkepov (\S\ref{sec:esirkepov}), and any spin-dipole curl current. Velocity
  is $\vv{v}=\vv{p}^{\,n-1/2}/(\gamma m)$. Optionally binomial-smooth
  $\rho,\vv{J}$ (\S\ref{sec:esirkepov}).
  \item \textbf{Add EM stress-energy} to the gravitational sources,
  $\rho_{\rm matter}\!\mathrel{+}=\rho_{\rm EM}$, $\vv{J}_{\rm matter}\!\mathrel{+}=\vv{J}_{\rm EM}$
  (Eqs.~\eqref{eq:rho_em}--\eqref{eq:J_em}).
  \item \textbf{Refresh $\clocal$} from the current total $\Phig$ if dynamic
  Shapiro is on (\S\ref{sec:em}).
  \item \textbf{Advance the fields:} one leapfrog sweep of all six potentials
  (\S\ref{sec:lattice}).
  \item \textbf{Absorb:} apply the centered-implicit boundary friction to each
  \texttt{next} array (\S\ref{sec:boundary}).
  \item \textbf{Rotate} the three time-level pointers of each field.
  \item \textbf{Gauge fixes:} optional zero-mean scalar subtraction and
  parabolic Lorenz cleaning (\S\ref{sec:boundary}).
  \item \textbf{Push particles:} gather the force, Boris kick--drift, then update
  spin, proper time, and position (\S\ref{sec:pusher}).
\end{algsteps}

The leapfrog staggering puts particle momenta/velocities at half-integer times
($\vv{p}^{\,n-1/2}$) and positions/fields at integer times; this is what makes
the deposition velocity, the force velocity, and the current-deposition segment
mutually consistent.

%====================================================================
\section{The Yee Lattice and the Leapfrog Field Update}
\label{sec:lattice}
%====================================================================

\subsection{Sublattices}

Each potential lives on its own staggered (Yee) sublattice so that the derived
fields land where they are needed:
\begin{center}
\begin{tabular}{@{}L{4.8cm}L{3.4cm}L{4.6cm}@{}}
\toprule
\textbf{Potential} & \textbf{Sublattice} & \textbf{Node position} \\
\midrule
$\Phig,\ \phi$       & corner  & $(i,\,j)\,\dx$ \\
$A_{g,x},\ A_x$      & $x$-edge & $(i{+}\tfrac12,\,j)\,\dx$ \\
$A_{g,y},\ A_y$      & $y$-edge & $(i,\,j{+}\tfrac12)\,\dx$ \\
\bottomrule
\end{tabular}
\end{center}
Curls then live at cell centers $(i{+}\tfrac12,j{+}\tfrac12)\dx$: e.g.\
$B_{g,z}=(A_{g,y}[i{+}1,j]-A_{g,y}[i,j])/\dx - (A_{g,x}[i,j{+}1]-A_{g,x}[i,j])/\dx$,
which the force gather evaluates on the fly (\S\ref{sec:gather}). The five-point
Laplacian at a field's own index sums neighbors on the \emph{same} sublattice,
so a single kernel form serves all six potentials; only the interpretation of
$(i,j)$ differs.

\subsection{The discrete update}

Writing $\Box\Psi = c^{-2}\partial_t^2\Psi-\nabla^2\Psi = \mathcal S$ for any of
the six potentials, the explicit leapfrog is
\begin{eqbox}
\begin{equation}
\label{eq:leapfrog}
  \Psi^{n+1}_{k} = 2\Psi^n_k - \Psi^{n-1}_k
     + (c\,\dt)^2\Bigl[(\nabla_h^2\Psi^n)_k + s_\Psi\,\mathcal{S}^n_k\Bigr],
  \quad
  (\nabla_h^2\Psi)_k = \frac{\Psi_{k\pm1}+\Psi_{k\pm W}-4\Psi_k}{\dx^2},
\end{equation}
\end{eqbox}
where $k=jW+i$ is the flat array index and $s_\Psi$ is the per-field source
coupling: $s_{\Phig}=-4\pi G$, $s_{\Ag}=-4\pi G/c^2$ (gravity), and
$s_{\phi}=1/\varepsilon_0$, $s_{\vv{A}}=-\mu_0$ (EM), realizing the source rows
of Table~\ref{tab:sixfield}. The scheme is second-order in space and time and
stable under the 2D Courant condition
\begin{equation}
\label{eq:cfl}
  \frac{c\,\dt}{\dx}\le \frac{1}{\sqrt2},
\end{equation}
exposed as the dimensionless \texttt{cfl} parameter ($\dt = \mathrm{cfl}\cdot
\dx/c$; the library default is $\mathrm{cfl}=1/\sqrt2$).

\subsection{Kernel variants}

The same recurrence \eqref{eq:leapfrog} appears in four forms, selected per run:
\begin{itemize}\setlength\itemsep{2pt}
  \item \textbf{undamped} (interior default): Eq.~\eqref{eq:leapfrog} verbatim;
  \item \textbf{Shapiro} (EM fields only): $c^2$ on the Laplacian is replaced by
  the per-cell $\clocal^2(\vv{x})$ array (\S\ref{sec:em}), while the source term
  keeps the uniform $c^2$;
  \item \textbf{critical damping}: with per-cell $\gamma_k=1-\sigma_k\dt$,
  $\Psi^{n+1}_k = 2\gamma_k\Psi^n_k-\gamma_k^2\Psi^{n-1}_k+(c\dt)^2[\nabla_h^2\Psi^n+s\mathcal S]$,
  whose double root at $\gamma_k$ gives an exact per-step decay and reduces
  bit-for-bit to the undamped form where $\sigma_k=0$;
  \item \textbf{periodic}: wrap-around indexing on all four edges (translation
  invariant).
\end{itemize}
For the non-periodic variants an outer-boundary post-pass sets the ring either to
zero (Dirichlet) or to its inward neighbor (Neumann); \S\ref{sec:boundary}
discusses the choice and the separate derivative-friction absorber.

%====================================================================
\section{Shape Functions and the Macroparticle}
\label{sec:shapes}
%====================================================================

A macroparticle replaces the point $\delta(\vv{x}-\vv{x}_p)$ of Part~I by a
finite cloud $W(\vv{x}-\vv{x}_p)$. All deposition and interpolation use a
separable shape $W(\vv{x})=w(x)\,w(y)$ whose 1D factor is non-negative,
symmetric, and a partition of unity, $\sum_i w(x-x_i)=1$ for any $x$ --- the
last property is what makes the deposited mass/charge independent of sub-cell
position. GRlite offers three shapes, all on the corner sublattice for $\rho$
and edge sublattices for $\vv J$:

\paragraph{CIC ($W_2$, bilinear, $2{\times}2$).} With
$\alpha=x_p/\dx-i_c$, $\beta=y_p/\dx-j_c$ the four weights are
$(1{-}\alpha)(1{-}\beta),\ \alpha(1{-}\beta),\ (1{-}\alpha)\beta,\ \alpha\beta$,
deposited as $\mathrm{value}\times w/\dx^2$ (so $\sum_k\!\rho_k\,\dx^2=\mathrm{value}$).

\paragraph{TSC ($W_3$, quadratic B-spline, $3{\times}3$).} Anchored at the
nearest node with $u=x_p/\dx-i_c\in[-\tfrac12,\tfrac12)$,
\begin{equation}
  w_{-1}=\tfrac12\bigl(\tfrac12-u\bigr)^2,\quad
  w_{0}=\tfrac34-u^2,\quad
  w_{+1}=\tfrac12\bigl(\tfrac12+u\bigr)^2 ,
\end{equation}
a smoother $9$-cell footprint with a continuous first derivative.

\paragraph{Bump ($C^\infty$, radius-parameterized).} A mollifier
\begin{equation}
\label{eq:bump}
  b(d,R)=\exp\!\Bigl(\frac{-1}{1-(d/R)^2}\Bigr)\ \text{for }|d|<R,\ \text{else }0,
  \qquad w_i = \frac{b(i-u,R)}{\sum_j b(j-u,R)},
\end{equation}
per-axis renormalized so the weights sum to one (hence charge-conserving). The
window half-width is $\lceil R+\tfrac12\rceil$ cells. Being $C^\infty$ with
compact support, its Fourier transform decays sub-exponentially, suppressing the
grid-scale aliasing that drives PIC self-heating; the production setup uses the
bump kernel at radius $R\approx6$. The same $w_i$ are reused for interpolation,
and the analytic derivative $\dd w_i/\dd u$ (including the renormalization
correction) is used for the gradient gather of \S\ref{sec:gather}.

\paragraph{The adjoint condition.} Deposit and gather use the \emph{same} shape;
this Hockney--Eastwood adjoint pairing is what guarantees a particle exerts no
net self-force in a uniform field, and underlies the momentum-conserving force
gather (\S\ref{sec:gather}).

%====================================================================
\section{Charge-Conserving Current Deposition}
\label{sec:esirkepov}
%====================================================================

\subsection{Why exact continuity matters}

The EM (and GEM) wave solver enforces no constraint that ties $\vv J$ to
$\partial_t\rho$. If the deposited pair violates the \emph{discrete} continuity
equation
\begin{equation}
\label{eq:discrete_continuity}
  \frac{\rho^{n}_k-\rho^{n-1}_k}{\dt} + (\nabla_h\!\cdot\vv J^{\,n-1/2})_k = 0,
\end{equation}
the excess charge sources a spurious longitudinal field that trails the particle
as a self-force wake. GRlite therefore deposits $\vv J$ by the Esirkepov
construction, which satisfies Eq.~\eqref{eq:discrete_continuity} \emph{exactly},
to round-off, at every corner cell.

\subsection{The Esirkepov split}

Let $S^{n-1},S^{n}$ be the shape footprints at $\vv{x}^{n-1},\vv{x}^{n}$, each a
product $S_x S_y$. The 2D footprint change
$\Delta S=S^{n}_xS^{n}_y-S^{n-1}_xS^{n-1}_y$
is split into orthogonal $x$- and $y$-flux parts
\begin{align}
  \Delta S^{x}_{ij} &= (S^{n}_{x,i}-S^{n-1}_{x,i})\bigl(S^{n-1}_{y,j}+\tfrac12(S^{n}_{y,j}-S^{n-1}_{y,j})\bigr),\\
  \Delta S^{y}_{ij} &= (S^{n}_{y,j}-S^{n-1}_{y,j})\bigl(S^{n-1}_{x,i}+\tfrac12(S^{n}_{x,i}-S^{n-1}_{x,i})\bigr),
\end{align}
which satisfy $\Delta S^x+\Delta S^y=\Delta S$ identically. The edge currents are
then the running sums
\begin{equation}
\label{eq:esirkepov_J}
  J_{x}\big|_{i+1/2,\,j} = -\frac{\mathrm{value}}{\dt\,\dx}\sum_{k\le i}\Delta S^x_{kj},
  \qquad
  J_{y}\big|_{i,\,j+1/2} = -\frac{\mathrm{value}}{\dt\,\dx}\sum_{l\le j}\Delta S^y_{il},
\end{equation}
with $\mathrm{value}=m$ (matter current) or $q$ (charge current). Telescoping the
sums reproduces Eq.~\eqref{eq:discrete_continuity} cell-by-cell. The trajectory
endpoints use $\vv{x}^{n-1}=\vv{x}^n-\vv{v}^{n-1/2}\dt$, exact under the leapfrog
drift. CIC ($2{\times}2$ footprint) and bump variants are provided; if the
Esirkepov support is exceeded (motion $>1$ cell), the code falls back to a direct
midpoint CIC current and flags the continuity violation.

\subsection{Spin-dipole and smoothing}

A spinning particle adds a divergence-free magnetization current. For a
$z$-dipole moment $\mu$ deposited with shape $W$,
\begin{equation}
  J_x = +\mu\,\partial_y W,\qquad J_y = -\mu\,\partial_x W
  \quad\Rightarrow\quad \nabla\!\cdot\vv J = 0,
\end{equation}
so it perturbs $\vv A$ without adding to $\rho$ (magnetic moment
$\mu=g_s q/2m\cdot S$ for EM, gravitomagnetic moment $\sigma=2S$). Finally, an
optional binomial filter --- the $3{\times}3$ stencil
$\tfrac1{16}[\,1\,2\,1;\,2\,4\,2;\,1\,2\,1\,]$, applied $N$ times to $\rho$ (and
to $\vv J$) --- damps the residual grid-scale content of the deposit, the
canonical PIC noise-control step. It is exposed as \texttt{rhoSmooth} /
\texttt{jSmooth} pass counts.

%====================================================================
\section{Force Gather and the Lewis--Birdsall Scheme}
\label{sec:gather}
%====================================================================

Because the state is potentials, the force \eqref{eq:force_gem_pot}--%
\eqref{eq:force_em_pot} is assembled at each particle from quantities gathered
on the fly --- no $\vv g,\Bg,\vv E,\vv B$ grid arrays are ever formed:
\begin{itemize}\setlength\itemsep{1pt}
  \item $\nabla\Phig$ and $\nabla\phi$: gathered gradients of the corner scalars;
  \item $\partial_t\Ag,\ \partial_t\vv A$: centered time differences of the
  perturbation edge potentials, $(\Psi^{n+1}-\Psi^{n-1})/2\dt$;
  \item $B_{g,z},B_z$: the Yee curl (\S\ref{sec:lattice}) evaluated at the four
  cell centers around the particle and bilinearly interpolated;
\end{itemize}
each combining the analytic (or sampled) background with the grid perturbation.

\subsection{Legacy vs.\ Lewis--Birdsall interpolation}

Two gather styles are selectable. The \emph{legacy} style interpolates field
values with the shape $W$ and finite-differences them. The
\emph{Lewis--Birdsall} (LB) style instead gathers the gradient \emph{directly}
with the analytic derivative of the same shape,
\begin{equation}
\label{eq:lb_grad}
  \partial_x\Psi(\vv{x}_p) = \frac1{\dx}\sum_{i,j}
     \frac{\dd w}{\dd u}(x_p-x_i)\,w(y_p-y_j)\,\Psi_{ij},
\end{equation}
and analogously for $\partial_y$ and for the edge-field gathers entering
$\vv{v}\times\vv B$. Pairing the derivative-shape gather \eqref{eq:lb_grad} with
the value-shape deposit is the adjoint (Hockney--Eastwood) condition at the
gradient level: it makes the inter-particle force exactly reciprocal and cancels
the grid self-force, so LB is the production choice. The 1PN gravitational
vector $\vv g=-\nabla\Phig-\partial_t\Ag$ used in the EIH terms is built from the
same gathered pieces.

\subsection{The assembled 2D force}

Collecting Part~I, the per-particle force the pusher applies is (relativistic
tier; $\vv g$ from the gathers above, $B_{g,z},B_z$ from the Yee curl)
\begin{eqbox}
\begin{align*}
  F_x &= m\Bigl[(1+\tfrac{v^2}{c^2}+\tfrac{4\Phig}{c^2})\,g_x
        - \tfrac{4(\vv v\cdot\vv g)}{c^2}v_x + 4 v_y B_{g,z} - \partial_tA_{g,x}\Bigr]
        + q\bigl[E_x + v_y B_z\bigr] + s\,\partial_x B_{g,z} + \mu\,\partial_x B_z,\\
  F_y &= m\Bigl[(1+\tfrac{v^2}{c^2}+\tfrac{4\Phig}{c^2})\,g_y
        - \tfrac{4(\vv v\cdot\vv g)}{c^2}v_y - 4 v_x B_{g,z} - \partial_tA_{g,y}\Bigr]
        + q\bigl[E_y - v_x B_z\bigr] + s\,\partial_y B_{g,z} + \mu\,\partial_y B_z,
\end{align*}
\end{eqbox}
with $E_\alpha=-\partial_\alpha\phi-\partial_tA_\alpha$, $\mu=g_s q/2m\cdot s$.
Lower tiers drop the $\tfrac{v^2}{c^2},\tfrac{4\Phig}{c^2},(\vv v\cdot\vv g)$
terms (EIH$\to$GEM) and the $v\times B_g$ term (GEM$\to$Newtonian); the
EM and spin pieces are unaffected. Each gravitomagnetic cross product is in
practice evaluated from the potential-form identity (\S\ref{sec:force_potential})
so $\Bg$ is never gridded.

%====================================================================
\section{The Relativistic Boris Pusher}
\label{sec:pusher}
%====================================================================

\subsection{Kick--drift with a velocity corrector}

Particles advance by a relativistic Boris leapfrog (kick--drift):
\begin{equation}
  \vv p^{\,n+1/2}=\vv p^{\,n-1/2}+\vv F^{\,n}\dt,\quad
  \vv v^{\,n+1/2}=\frac{\vv p^{\,n+1/2}}{\sqrt{m^2+|\vv p^{\,n+1/2}|^2/c^2}},\quad
  \vv x^{\,n+1}=\vv x^{\,n}+\vv v^{\,n+1/2}\dt.
\end{equation}
The force depends on velocity (the $v^2/c^2$ EIH terms and the $\vv v\times\vv B$
terms), which is not known at $t^n$. GRlite uses the lagged half-step velocity
$\vv v^{\,n-1/2}$ as a predictor, then, whenever the force is velocity-dependent
(relativistic tier, or $B_{g,z}\neq0$, or a charged particle in $B_z$), runs one
corrector: predict $\vv p$, form the midpoint $\vv v_{\rm mid}=\tfrac12(\vv
v^{\,n-1/2}+\vv v_{\rm pred})$, and re-evaluate the force there. This restores
second-order time accuracy and the time-symmetry of the push.

\subsection{Self-field subtraction}

For particles that opt in (\S\ref{sec:bgimpl}), the gathered force includes the
particle's own deposited field --- a residual grid self-force. GRlite cancels it
by maintaining a per-particle ``self'' field set (the same deposit and leapfrog,
sourced by that particle alone) and forming
\begin{equation}
\label{eq:selffield}
  \vv F_{\rm total} = \vv F_{\rm collective} - (1+\varepsilon)\,\vv F_{\rm self},
\end{equation}
re-gathered through the identical chain; $\varepsilon=0$ gives exact bit-level
cancellation of the self-interaction, leaving only the inter-particle force.

\subsection{Spin, proper time, and driven sources}

After the push, the scalar spin phase advances by the 2D precession rate
(Eq.~\eqref{eq:precession_rates}, $\Omega_z=B_{g,z}+\tfrac{g_sq}{2m}B_z+
\Omega_{{\rm Thomas},z}$ with $\Omega_{{\rm Thomas},z}=-\tfrac{\gamma^2}{\gamma+1}
(v_xa_y-v_ya_x)/c^2$ and $\vv a=(\vv v^{\,n+1/2}-\vv v^{\,n-1/2})/\dt$), and the
proper-time clock by
\begin{equation}
  \Delta\tau = \dt\sqrt{1+\tfrac{2\Phig}{c^2}-\bigl(1-\tfrac{2\Phig}{c^2}\bigr)\tfrac{v^2}{c^2}-\tfrac{8\,\Ag\cdot\vv v}{c^2}}
\end{equation}
(Eq.~\eqref{eq:proper_time}). Two source modes override the plain push: a
\emph{pinned} particle (forces disabled) does not integrate the force, acting as
a fixed emitter; a \emph{driven} particle carries a base velocity $\vv v_{\rm
base}$ (itself evolved under the force) plus a sinusoidal modulation $\vv v=\vv
v_{\rm base}+\vv v_{\rm drive}\cos(\omega t+\varphi)$, written into the momentum
so the Esirkepov current stays charge-conserving. Storing the base separately is
what prevents the oscillation from accumulating into the momentum --- the model
for an antenna/dipole emitter.

%====================================================================
\section{Initial Conditions}
\label{sec:init}
%====================================================================

Starting the leapfrog from zero fields with sources already present injects a
spurious ``turn-on'' transient: the field has to radiate out to its static
configuration, and the standing modes of a finite box are slow to leave. GRlite
seeds a near-static initial field instead, by two ingredients:

\paragraph{Liénard--Wiechert seed.} The potentials are initialized to the
direct-sum (here static) Liénard--Wiechert superposition of the particle
sources, with a boundary-mean taper that matches the interior solution smoothly
to the box edge (matched to zero across the damping ring by default). This places
the field within a small residual of the discrete static solution at $t=0$.

\paragraph{Frozen-particle friction settle.} The L--W seed is then relaxed to the
exact discrete fixed point by running $N_{\rm settle}$ steps with the particles
\emph{frozen} (sources still deposit from the frozen state; the field still
evolves) and the derivative-friction absorber (\S\ref{sec:boundary}) on. The
friction drains the residual standing modes to round-off in a few hundred steps,
leaving the field at the discrete Poisson solution of the deposited $\rho$.

These compose into the scenario-level initialization choices
\texttt{lw-settled} (seed $+$ settle, the default for interacting runs),
\texttt{lw} (seed only), \texttt{none} (auto: settle when a perturbation is
active, else leave at zero), and \texttt{zero} (always start from a quiet field
--- used by wave-emission/antenna demos that must not begin with a static well).
A direct Jacobi/Poisson relaxation of $\Phig$ is also available as an
alternative seed.

%====================================================================
\section{Boundary Conditions and Gauge Control}
\label{sec:boundary}
%====================================================================

\subsection{Outer boundary}

The non-periodic solver applies, after each leapfrog sweep, either a
\textbf{Dirichlet} ring ($\Psi=0$ on the outer cells: waves reflect with a sign
flip; the discrete static state is the zero-Dirichlet Poisson solution) or a
\textbf{Neumann} ring ($\Psi_{\rm edge}=\Psi_{\rm inward}$: reflection without
sign flip, putting standing-mode antinodes at the wall where the absorber is
strongest, and better approximating the 2D free-space $\ln r$ Green's function).
A fully \textbf{periodic} variant (wrap-around indexing) is also available.

\subsection{The absorbing layer}

Outgoing waves are removed by a derivative-friction layer implementing the damped
wave equation $\partial_t^2\Psi+2\gamma\,\partial_t\Psi=c^2\nabla^2\Psi+\dots$.
The discretization is \emph{centered-implicit}: the standard leapfrog
\texttt{next} is corrected, before the buffer rotation, by
\begin{equation}
\label{eq:friction}
  \Psi^{n+1}_k \leftarrow \frac{\Psi^{n+1}_k + \beta_k\,\Psi^{n-1}_k}{1+\beta_k},
  \qquad \beta_k=\gamma_k\,\dt .
\end{equation}
At the 2D Nyquist mode the recurrence roots are $-1$ (marginal) and
$-(1-\beta)/(1+\beta)$ (stable), so \eqref{eq:friction} is stable for \emph{any}
$\beta$ --- unlike a naive bleed on \texttt{prev}, which destabilizes the Nyquist
mode at $\mathrm{cfl}=1/\sqrt2$. The coefficient $\beta_k$ is a cubic ramp from a
small interior \emph{floor} (which also drains the lowest box mode) up to a
\emph{wall} value over a taper \emph{depth} of cells at the boundary.

\subsection{Gauge maintenance}

Two optional fixes keep the Lorenz gauge from drifting. The \textbf{zero-mean}
fix addresses the unconstrained DC mode of the scalar potentials under Neumann
BC (the wave equation accelerates it whenever the source mean is nonzero):
subtract the spatial mean from both \texttt{curr} and \texttt{prev} each step
(equal subtraction preserves the encoded time derivative). The
\textbf{parabolic Lorenz cleaning} subtracts a multiple of the gauge residual,
$\phi\mathrel{-}=R\,c^2\dt\,(\nabla\!\cdot\vv A)$ (and likewise $\Phig$ with
$\nabla\!\cdot\Ag$), bounding the drift introduced by sources whose discrete
divergence is not exactly consistent (e.g.\ the spin-dipole curl current).

%====================================================================
\section{Background Field Installation}
\label{sec:bgimpl}
%====================================================================

The fixed background of \S\ref{sec:bgsplit} is installed from a list of up to
sixteen compact bodies, each with parameters $(x,y,GM,Q,J_z,g_s,\varepsilon)$.
The list is the source of truth; the field is derived from it in one of two
modes.

\paragraph{Sampled mode.} Each body's analytic generators
(\S\ref{sec:compactbody}) --- the softened gravitational scalar, the
gravitomagnetic dipole, the Coulomb scalar, and (for a spinning charge) the EM
magnetic dipole --- are sampled \emph{additively} onto the six background grid
arrays, on each potential's own sublattice. Multiple bodies superpose by
accumulation, which is exact in the linear theory. The force gather then reads
these grids exactly as it reads the perturbation, and the Shapiro $\clocal$
(Eq.~\eqref{eq:clocal}) is built from the \emph{total} $\Phig$ (sampled
background $+$ perturbation), so a deposited mass lenses light on top of the
fixed body.

\paragraph{Analytic mode.} The closed-form fields are instead evaluated directly
at each particle position and summed over the bodies, with no grid-sampling
error. This eliminates the $\order{(\dx/r)^2}$ tangential-force error of the
sampled path that would otherwise induce a slow grid precession --- so analytic
mode is preferred for clean test-particle orbits, while sampled mode is used when
the background must enter the gridded field evolution (e.g.\ dynamic lensing).
Either way, a single body reduces exactly to the corresponding classical-metric
analogue of \S\ref{sec:compactbody}.

%====================================================================
\section{Visualization Mathematics}
\label{sec:viz}
%====================================================================

The visualization is part of the physics: a pedagogical and debugging tool only
works if the rendered quantity is exactly the simulated one. All field views are
computed CPU-side each frame from the potential arrays and the chosen
\emph{source} (perturbation only, background only, or their sum).

\paragraph{Field colormap.} A selectable scalar is mapped through a diverging
colormap auto-normalized to the current frame's peak magnitude (so the scale
tracks the dynamics). The available reductions are the stored scalars themselves
($\Phig$, $\phi$); a \emph{curl} field
$B_z=\partial_xA_y-\partial_yA_x$ (and $B_{g,z}$) formed from the two edge
arrays; and a \emph{gradient magnitude} $|\nabla\Phig|$ for the field strength.

\paragraph{Vector overlay (quiver).} On a coarse sub-grid the overlay draws
either edge-potential components ($\Ag,\vv A$) or a gradient field
($\vv g=-\nabla\Phig$, $\vv E=-\nabla\phi$), arrows normalized so the strongest
spans $\approx0.9$ of the sample spacing, with opacity scaling by relative
magnitude so weak regions recede.

\paragraph{Per-particle overlays.} Each maps a simulated quantity to a glyph:
\begin{itemize}\setlength\itemsep{2pt}
  \item \textbf{velocity arrow} of pixel length $L_v\,|\vv v|$ with
  $\vv v=\vv p/(\gamma m)$ (a fixed pixels-per-$c$ scale), brightened above
  $|\vv v|>0.5c$;
  \item \textbf{force arrows} for the gravitational, EM, spin, and total
  components (the per-particle breakdown stashed by the pusher), drawn in
  distinct colors and sharing \emph{one} auto-scale (the largest component across
  all particles spans a fixed reference length) so magnitudes are directly
  comparable;
  \item \textbf{spin indicator}: a hand through the particle at angle
  $\phi_{\rm spin}$, so Lense--Thirring/Thomas precession is literally visible;
  \item \textbf{trajectory ticks}: along the recorded track, circles at equal
  \emph{coordinate}-time intervals $\Delta$ and triangles at equal
  \emph{proper}-time intervals $\Delta$ (auto-chosen to $\sim12$ ticks across the
  trail). Where the proper-time triangles spread apart relative to the
  coordinate circles, time dilation is being shown directly;
  \item \textbf{proper-time clock}: a two-hand dial, one hand sweeping coordinate
  time $t$ and one proper time $\tau$ at a common period; the opening wedge
  between them is the accumulated dilation $t-\tau$;
  \item \textbf{background body}: a hollow ring sized to the softening length
  $\varepsilon$ at the body's position.
\end{itemize}
Every overlay is independently toggleable so a single effect can be isolated.

\bigskip
\noindent\textit{End of Part~II.}  Part~III documents the C library that
implements the scheme above and maps each public function to the pipeline step
(\S\ref{sec:pipeline}) it realizes.

\part{The C Library}

The simulation core is a freestanding C11 library (\texttt{core/}, compiled both
natively for the test suite and to WebAssembly for the browser). It has no
dependencies beyond \texttt{libm}. This Part documents the data structures that
hold the state of Part~II, the source files that implement each piece of the
scheme, and the complete public API with each function tied back to the
algorithm step it realizes. The single public header is
\texttt{core/include/grlite.h}; the internal state struct lives in
\texttt{core/src/sim\_internal.h}.

%====================================================================
\section{Architecture and Data Structures}
\label{sec:arch}
%====================================================================

\subsection{The simulation object}

A run is a single opaque handle \texttt{gr\_sim\_t*}, created with the grid and
the effective constants of Part~I:
\begin{verbatim}
gr_sim_t* gr_sim_create(int width, int height, float dx, float c_eff, float cfl);
\end{verbatim}
Internally (\texttt{struct gr\_sim}) it holds the entire state of
\S\ref{sec:pipeline}:
\begin{center}
\begin{tabular}{@{}L{5.6cm}L{8.2cm}@{}}
\toprule
\textbf{Member} & \textbf{Role (Part~II reference)} \\
\midrule
\texttt{c\_eff, dt, cfl, dx} & wave speed / step / Courant ratio (Eq.~\eqref{eq:cfl}) \\
\texttt{G\_eff, k\_e}        & effective couplings $G,\keff$ (Part~I) \\
\texttt{fields[6]}           & the six potentials, each a \texttt{gr\_field\_state\_t} (\S\ref{sec:lattice}) \\
\texttt{rho\_matter, J\_mx, J\_my} & gravity sources $\rho_{\rm matter},\vv J_{\rm matter}$ (\S\ref{sec:esirkepov}) \\
\texttt{rho\_q, J\_qx, J\_qy}      & EM sources $\rho_q,\vv J_q$ \\
\texttt{particles, n\_particles}  & the macroparticle array (\S\ref{sec:pusher}) \\
\texttt{damping\_d, friction\_d}  & absorbing-layer arrays (\S\ref{sec:boundary}) \\
\texttt{phi\_g\_bg, \dots, Ay\_bg} & sampled background fields (\S\ref{sec:bgimpl}) \\
\texttt{bg\_bodies[16], bg\_nbody} & analytic compact-body list (\S\ref{sec:bgimpl}) \\
\texttt{self\_field\_sets}         & per-particle self-field sets (Eq.~\eqref{eq:selffield}) \\
\texttt{force\_tier, shape\_function, force\_interp} & scheme selectors (\S\S\ref{sec:shapes},\ref{sec:gather}) \\
\texttt{*\_enabled flags}          & the physics switches (force terms, Shapiro, gauge) \\
\bottomrule
\end{tabular}
\end{center}

\subsection{Field, particle, and body records}

Each potential is a three-time-level record on its Yee sublattice; the leapfrog
\eqref{eq:leapfrog} reads \texttt{prev}/\texttt{curr}, writes \texttt{next}, and
the step rotates the three pointers:
\begin{verbatim}
typedef struct { float *prev, *curr, *next;     /* Phi^{n-1}, Phi^n, Phi^{n+1} */
                 float *source; float source_coeff; } gr_field_state_t;
\end{verbatim}
\texttt{source} binds to one of the six source arrays and \texttt{source\_coeff}
is the $s_\Psi$ of Eq.~\eqref{eq:leapfrog}. A macroparticle carries the dynamical
state of \S\ref{sec:pusher} plus per-step diagnostics:
\begin{verbatim}
typedef struct { float x, y, px, py;            /* position, relativistic momentum */
                 float mass, charge, proper_time;
                 float spin, phi_spin, g_factor; /* scalar spin (Sec. 4.3) */
                 float fgrav_x..ftot_y;          /* force breakdown (viz, Sec. 16) */
                 float drive_vx..drive_base_vy;  /* driven/pinned source (Sec. 12) */
               } gr_particle_t;
\end{verbatim}
and a background body is the seven contiguous floats of the no-hair source
(\S\ref{sec:compactbody}):
\begin{verbatim}
typedef struct { float x, y, GM, Q, Jz, g_em, eps; } gr_bg_body_t;
\end{verbatim}

\subsection{Memory layout and access}

Every grid array is a flat \texttt{float[W*H]} in row-major order (index
$k=jW+i$), matching the leapfrog indexing of Eq.~\eqref{eq:leapfrog}. All grids
share one geometry; only the node interpretation (corner / edge) differs by
sublattice. Because the arrays and the particle records are plain
\texttt{float}s in a single linear heap, the WebAssembly build exposes them to
the front end as zero-copy \texttt{Float32Array} views: \texttt{gr\_sim\_field\_ptr},
\texttt{gr\_sim\_array\_ptr}, \texttt{gr\_sim\_background\_ptr}, and
\texttt{gr\_sim\_get\_particle} hand back interior pointers that the visualization
(\S\ref{sec:viz}) reads directly each frame.

\subsection{Lifecycle}

The driver loop is just \texttt{gr\_sim\_create} $\to$ configure $\to$ add
particles/background $\to$ initialize fields (\S\ref{sec:init}) $\to$
\texttt{gr\_sim\_step}/\texttt{gr\_sim\_step\_n} $\to$ \texttt{gr\_sim\_destroy}.
The convenience entry \texttt{gr\_sim\_use\_pedagogical\_defaults} installs the
production bundle (bump kernel, Lewis--Birdsall, deposition on, absorbing layer)
in one call.

%====================================================================
\section{Module-by-Module Correspondence}
\label{sec:modules}
%====================================================================

The implementation is six source files, each owning one band of Part~II
(Table~\ref{tab:modules}).

\begin{table}[h]
\centering
\begin{tabular}{@{}L{2.4cm}F{6.4cm}F{4.6cm}@{}}
\toprule
\textbf{File} & \textbf{Implements (Part~II)} & \textbf{Key symbols} \\
\midrule
\texttt{sim.c} & lifecycle; the per-step driver (\S\ref{sec:pipeline}); init (\S\ref{sec:init}); friction taper, zero-mean (\S\ref{sec:boundary}); config & \texttt{gr\_sim\_step}, \texttt{\_create}, \texttt{\_init\_potentials\_settled}, \texttt{\_set\_volume\_friction\_taper} \\
\texttt{field.c} & leapfrog + variants (\S\ref{sec:lattice}); parabolic gauge clean (\S\ref{sec:boundary}) & \texttt{gr\_field\_leapfrog\_step\_all}, \texttt{\_step\_self}, \texttt{gr\_field\_parabolic\_gauge\_clean} \\
\texttt{deposit.c} & shapes (\S\ref{sec:shapes}); Esirkepov (\S\ref{sec:esirkepov}); LB gathers (\S\ref{sec:gather}); spin-dipole curl & \texttt{gr\_\{cic,tsc,bump\}\_deposit\_corner}, \texttt{gr\_\{,bump\_\}esirkepov\_deposit\_jxy}, \texttt{gr\_bump\_lb\_grad\_corner} \\
\texttt{particle.c} & force gather (\S\ref{sec:gather}); Boris push, corrector, self-field, spin, proper time, drive (\S\ref{sec:pusher}) & \texttt{gr\_particle\_push\_all}, \texttt{gr\_phi\_g\_total\_at} \\
\texttt{background.c} & compact-body generators + analytic evaluators (\S\ref{sec:bgimpl}); Shapiro $\clocal$ rebuild (\S\ref{sec:em}) & \texttt{gr\_sim\_\{set,add\}\_background\_body}, \texttt{gr\_bg\_eval\_*}, \texttt{gr\_em\_shapiro\_recompute\_c\_local2} \\
\texttt{gauge.c} & Lorenz-gauge residual diagnostics (\S\ref{sec:boundary}) & \texttt{gr\_sim\_gauge\_residual\_\{grav,em\}} \\
\bottomrule
\end{tabular}
\caption{Source-file to Part~II mapping.}
\label{tab:modules}
\end{table}

Two correspondences are worth spelling out because they thread several files.

\paragraph{The step driver (\texttt{sim.c} $\Rightarrow$ \S\ref{sec:pipeline}).}
\texttt{gr\_sim\_step} executes the pipeline literally in order: it clears the
source arrays and, per particle, calls the selected
\texttt{gr\_*\_deposit\_corner} for $\rho^n$ and \texttt{gr\_*\_esirkepov\_deposit\_jxy}
for $\vv J^{\,n-1/2}$ (with \texttt{gr\_*\_curl\_dipole\_deposit\_jxy} for spin),
optionally binomial-smooths them, folds in EM stress-energy, refreshes $\clocal$
(\texttt{background.c}), calls \texttt{gr\_field\_leapfrog\_step\_all}
(\texttt{field.c}), applies the centered-implicit friction
(Eq.~\eqref{eq:friction}) to each \texttt{next} before the pointer rotation, runs
the gauge fixes, and finally calls \texttt{gr\_particle\_push\_all}
(\texttt{particle.c}) unless the particles are frozen.

\paragraph{The push (\texttt{particle.c} $\Rightarrow$ \S\S\ref{sec:gather}--\ref{sec:pusher}).}
\texttt{gr\_particle\_push\_all} gathers each force component through the static
evaluators (\texttt{grav\_grad\_at}/\texttt{\_lb}, \texttt{B\_g\_z\_at\_total},
\texttt{phi\_em\_grad\_at\_total}, \texttt{dt\_A\_*\_at\_total}), assembles the
tiered force, runs the velocity corrector, performs the self-field subtraction by
temporarily swapping \texttt{sim->fields} for the particle's
\texttt{self\_field\_set} and replaying the same gather, then updates momentum,
spin phase, proper time, and position. The force-component fields it writes back
(\texttt{fgrav/fem/fspin/ftot}) are exactly what the force-arrow overlay
(\S\ref{sec:viz}) renders.

%====================================================================
\section{API Reference}
\label{sec:api}
%====================================================================

The public API (\texttt{grlite.h}) is listed below by function group, each entry
tied to the Part~II step it serves. Arguments are abbreviated; \texttt{s}
denotes the \texttt{gr\_sim\_t*} handle. Getters that merely return the value set
by their paired setter are folded into one row (``\,set/get\,'').

\subsection{Lifecycle and time}
\nopagebreak
{\footnotesize
\begin{tabular}{@{}F{8.2cm}L{4.9cm}L{1.4cm}@{}}
\toprule
\textbf{Function} & \textbf{Role} & \textbf{Part~II} \\
\midrule
\texttt{gr\_sim\_create(w,h,dx,c\_eff,cfl)} & allocate a run & \S\ref{sec:arch} \\
\texttt{gr\_sim\_destroy(s)} & free it & \S\ref{sec:arch} \\
\texttt{gr\_sim\_use\_pedagogical\_defaults(s)} & install production bundle & \S\ref{sec:arch} \\
\texttt{gr\_sim\_step(s)} / \texttt{gr\_sim\_step\_n(s,n)} & advance the pipeline & \S\ref{sec:pipeline} \\
\texttt{gr\_sim\_step\_count(s)}, \texttt{gr\_sim\_time(s)} & step index, $t=n\,\dt$ & \S\ref{sec:pipeline} \\
\texttt{gr\_sim\_reset\_time(s)} & zero the step counter & \S\ref{sec:pipeline} \\
\texttt{gr\_sim\_dt/dx/width/height(s)} & grid/time accessors & \S\ref{sec:lattice} \\
\bottomrule
\end{tabular}}

\subsection{Couplings and scheme selectors}
\nopagebreak
{\footnotesize
\begin{tabular}{@{}F{8.2cm}L{4.9cm}L{1.4cm}@{}}
\toprule
\textbf{Function} & \textbf{Role} & \textbf{Part~II} \\
\midrule
\texttt{gr\_sim\_set/get\_G\_eff(s,\,$\cdot$)} & gravity coupling $G$ & Part~I \\
\texttt{gr\_sim\_set/get\_k\_e(s,\,$\cdot$)} & EM coupling $\keff$ & Part~I \\
\texttt{gr\_sim\_set/get\_force\_tier(s,\,$\cdot$)} & Newtonian / GEM / EIH & \S\ref{sec:gather} \\
\texttt{gr\_sim\_set/get\_shape\_function(s,\,$\cdot$)} & CIC / TSC / bump & \S\ref{sec:shapes} \\
\texttt{gr\_sim\_set/get\_kernel\_radius(s,\,$R$)} & bump radius $R$ & \S\ref{sec:shapes} \\
\texttt{gr\_sim\_set/get\_force\_interp(s,\,$\cdot$)} & legacy / Lewis--Birdsall & \S\ref{sec:gather} \\
\bottomrule
\end{tabular}}

\subsection{Physics switches}
Each toggles one term of the assembled force, source, or gauge machinery; all are
\texttt{set/get} pairs taking \texttt{(s, enabled)}.

\nopagebreak
{\footnotesize
\begin{tabular}{@{}F{8.2cm}L{4.9cm}L{1.4cm}@{}}
\toprule
\textbf{Function (\,\texttt{gr\_sim\_set\_}$\cdots$\texttt{\_enabled}\,)} & \textbf{Gates} & \textbf{Part~II} \\
\midrule
\texttt{gravitomagnetic\_force} & $4m\,\vv v\times\Bg$ & \S\ref{sec:gather} \\
\texttt{gravitomagnetic\_inductive} & $-m\,\partial_t\Ag$ & \S\ref{sec:gather} \\
\texttt{em\_lorentz\_force} & $q(\vv E+\vv v\times\vv B)$ & \S\ref{sec:gather} \\
\texttt{em\_inductive} & $-q\,\partial_t\vv A$ & \S\ref{sec:gather} \\
\texttt{em\_electrostatic} & $-q\nabla\phi$ & \S\ref{sec:gather} \\
\texttt{em\_magnetic} & $q\,\vv v\times\vv B$ & \S\ref{sec:gather} \\
\texttt{em\_stress\_energy} & $\rho_{\rm EM},\vv J_{\rm EM}$ source & \S\ref{sec:pipeline} \\
\texttt{em\_shapiro} & $\clocal$ EM solve & \S\ref{sec:em} \\
\texttt{em\_shapiro\_dynamic} & per-step $\clocal$ rebuild & \S\ref{sec:bgimpl} \\
\texttt{parabolic\_gauge\_cleaning} & Lorenz cleaning & \S\ref{sec:boundary} \\
\texttt{field\_evolution} & run the leapfrog at all & \S\ref{sec:lattice} \\
\texttt{particle\_source\_deposition} & deposit particle sources & \S\ref{sec:esirkepov} \\
\texttt{esirkepov} & charge-conserving $\vv J$ & \S\ref{sec:esirkepov} \\
\texttt{periodic\_bc} & wrap-around boundary & \S\ref{sec:boundary} \\
\texttt{particles\_frozen} & freeze the push (settle) & \S\ref{sec:init} \\
\bottomrule
\end{tabular}}

\subsection{Sources and deposition}
\nopagebreak
{\footnotesize
\begin{tabular}{@{}F{8.2cm}L{4.9cm}L{1.4cm}@{}}
\toprule
\textbf{Function} & \textbf{Role} & \textbf{Part~II} \\
\midrule
\texttt{gr\_sim\_clear\_sources(s)} & zero the six source arrays & \S\ref{sec:pipeline} \\
\texttt{gr\_sim\_clear\_fields(s)} & zero all potentials & \S\ref{sec:init} \\
\texttt{gr\_sim\_deposit\_point\_mass/charge(s,x,y,$\cdot$)} & direct $\rho$ deposit (shape) & \S\ref{sec:shapes} \\
\texttt{gr\_sim\_deposit\_point\_particle(s,\dots)} & deposit $\rho$ and $\vv J$ & \S\ref{sec:esirkepov} \\
\texttt{gr\_sim\_rho\_matter\_ptr/rho\_q\_ptr(s)} & source-array views & \S\ref{sec:esirkepov} \\
\texttt{gr\_sim\_J\_mx/my/qx/qy\_ptr(s)} & current-array views & \S\ref{sec:esirkepov} \\
\texttt{gr\_sim\_set/get\_rho\_smooth\_passes(s,n)} & $\rho$ binomial smoothing & \S\ref{sec:esirkepov} \\
\texttt{gr\_sim\_set/get\_j\_smooth\_passes(s,n)} & $\vv J$ binomial smoothing & \S\ref{sec:esirkepov} \\
\texttt{gr\_sim\_esirkepov\_violations(s)} & continuity-fallback count & \S\ref{sec:esirkepov} \\
\bottomrule
\end{tabular}}

\subsection{Field / array access and gauge diagnostics}
\nopagebreak
{\footnotesize
\begin{tabular}{@{}F{8.2cm}L{4.9cm}L{1.4cm}@{}}
\toprule
\textbf{Function} & \textbf{Role} & \textbf{Part~II} \\
\midrule
\texttt{gr\_sim\_field\_ptr(s,which)} & a potential array (\texttt{prev/curr}) & \S\ref{sec:lattice} \\
\texttt{gr\_sim\_array\_ptr(s,which)} & any field/source array & \S\ref{sec:arch} \\
\texttt{gr\_sim\_background\_ptr(s,which)} & a sampled background array & \S\ref{sec:bgimpl} \\
\texttt{gr\_array\_lattice(which)} & sublattice of an array & \S\ref{sec:lattice} \\
\texttt{gr\_lattice\_offset(lat,$\cdot$,$\cdot$)} & node half-cell offset & \S\ref{sec:lattice} \\
\texttt{gr\_sim\_gauge\_residual\_grav/em(s)} & $\|\partial_\mu\bar h^{\mu\nu}\|$ monitor & \S\ref{sec:boundary} \\
\bottomrule
\end{tabular}}

\subsection{Particles}
\nopagebreak
{\footnotesize
\begin{tabular}{@{}F{8.2cm}L{4.9cm}L{1.4cm}@{}}
\toprule
\textbf{Function} & \textbf{Role} & \textbf{Part~II} \\
\midrule
\texttt{gr\_sim\_add\_particle(s,x,y,m,q,vx,vy)} & append a macroparticle & \S\ref{sec:pusher} \\
\texttt{gr\_sim\_particle\_count(s)} & number of particles & \S\ref{sec:pusher} \\
\texttt{gr\_sim\_get\_particle(s,idx)} & read a particle record & \S\ref{sec:viz} \\
\texttt{gr\_sim\_clear\_particles(s)} & remove all & \S\ref{sec:pusher} \\
\texttt{gr\_sim\_particle\_energy(s,idx)} & $\gamma m c^2$ diagnostic & \S\ref{sec:pusher} \\
\texttt{gr\_sim\_set\_particle\_spin(s,idx,s,g)} & scalar spin + $g_s$ & \S\ref{sec:spin} \\
\texttt{gr\_sim\_set\_particle\_phi\_spin(s,idx,$\phi$)} & spin phase & \S\ref{sec:spin} \\
\texttt{gr\_sim\_set\_particle\_drive(s,idx,\dots)} & driven/pinned source & \S\ref{sec:pusher} \\
\texttt{gr\_sim\_particle\_enable/disable\_self\_field(s,idx)} & toggle self-field set & Eq.~\eqref{eq:selffield} \\
\texttt{gr\_sim\_particle\_has\_self\_field(s,idx)} & query & Eq.~\eqref{eq:selffield} \\
\texttt{gr\_sim\_particle\_set/get\_self\_field\_epsilon(s,idx,$\cdot$)} & subtraction $\varepsilon$ & Eq.~\eqref{eq:selffield} \\
\bottomrule
\end{tabular}}

\subsection{Initialization}
\nopagebreak
{\footnotesize
\begin{tabular}{@{}F{8.2cm}L{4.9cm}L{1.4cm}@{}}
\toprule
\textbf{Function} & \textbf{Role} & \textbf{Part~II} \\
\midrule
\texttt{gr\_sim\_relax\_phi\_g\_poisson(s,n)} & Jacobi Poisson seed & \S\ref{sec:init} \\
\texttt{gr\_sim\_init\_potentials\_lienard\_wiechert(s)} & L--W seed & \S\ref{sec:init} \\
\texttt{gr\_sim\_init\_potentials\_lienard\_wiechert\_with\_taper(s,in,out)} & L--W seed + boundary taper & \S\ref{sec:init} \\
\texttt{gr\_sim\_init\_potentials\_settled(s,n\_settle)} & L--W + frozen friction settle & \S\ref{sec:init} \\
\bottomrule
\end{tabular}}

\subsection{Boundary, absorber, and damping}
\nopagebreak
{\footnotesize
\begin{tabular}{@{}F{8.2cm}L{4.9cm}L{1.4cm}@{}}
\toprule
\textbf{Function} & \textbf{Role} & \textbf{Part~II} \\
\midrule
\texttt{gr\_sim\_set/get\_outer\_bc\_neumann(s,$\cdot$)} & Dirichlet vs Neumann ring & \S\ref{sec:boundary} \\
\texttt{gr\_sim\_set\_volume\_friction\_taper(s,floor,wall,depth)} & derivative-friction absorber & \S\ref{sec:boundary} \\
\texttt{gr\_sim\_volume\_friction\_enabled(s)} & query & \S\ref{sec:boundary} \\
\texttt{gr\_sim\_set/get\_zero\_mean\_scalar\_potentials(s,$\cdot$)} & DC-mode gauge fix & \S\ref{sec:boundary} \\
\texttt{gr\_sim\_set\_damping(s,n)} / \texttt{\_damping\_layers(s)} & PML width & \S\ref{sec:boundary} \\
\texttt{gr\_sim\_set/get\_damping\_config(s,cfg)} & PML profile params & \S\ref{sec:boundary} \\
\texttt{gr\_sim\_damping\_max\_stable\_sigma\_dt(s,form)} & stability bound & \S\ref{sec:boundary} \\
\bottomrule
\end{tabular}}

\subsection{Background fields}
The unified compact body and the multi-body calls implement \S\ref{sec:bgimpl};
the uniform/linear setters are single-body isolation fixtures.

\nopagebreak
{\footnotesize
\begin{tabular}{@{}F{8.2cm}L{4.9cm}L{1.4cm}@{}}
\toprule
\textbf{Function} & \textbf{Role} & \textbf{Part~II} \\
\midrule
\texttt{gr\_sim\_clear\_background(s)} & drop all background & \S\ref{sec:bgimpl} \\
\texttt{gr\_sim\_set\_background\_body(s,x,y,GM,Q,Jz,g,eps)} & one $(M,Q,J)$ body & \S\ref{sec:compactbody} \\
\texttt{gr\_sim\_add\_background\_body(s,\dots)} $\to$ idx & append a body (superpose) & \S\ref{sec:bgimpl} \\
\texttt{gr\_sim\_background\_count(s)} & number of bodies & \S\ref{sec:bgimpl} \\
\texttt{gr\_sim\_get\_background\_body\_ptr(s,i)} & body params view (7 floats) & \S\ref{sec:bgimpl} \\
\texttt{gr\_sim\_set\_background\_body\_at(s,i,\dots)} & edit body $i$ (re-derive) & \S\ref{sec:bgimpl} \\
\texttt{gr\_sim\_set\_background\_point\_mass/\_charge(s,\dots)} & Schwarzschild / RN special case & \S\ref{sec:compactbody} \\
\texttt{gr\_sim\_set\_background\_spinning\_point\_mass(s,\dots)} & Kerr special case & \S\ref{sec:compactbody} \\
\texttt{gr\_sim\_set\_background\_uniform\_gravitomagnetic/\_magnetic/\_electric(s,\dots)} & uniform test fields & \S\ref{sec:bgimpl} \\
\texttt{gr\_sim\_set\_background\_linear\_magnetic(s,\dots)} & $B_z(x)$ gradient fixture & \S\ref{sec:bgimpl} \\
\texttt{gr\_sim\_set/get\_bg\_mode(s,$\cdot$)} & sampled vs analytic & \S\ref{sec:bgimpl} \\
\texttt{gr\_sim\_get\_bg\_kind(s)} & installed background kind & \S\ref{sec:bgimpl} \\
\bottomrule
\end{tabular}}

\subsection{Scenario registry}
\nopagebreak
{\footnotesize
\begin{tabular}{@{}F{8.2cm}L{4.9cm}L{1.4cm}@{}}
\toprule
\textbf{Function} & \textbf{Role} & \textbf{Part~II} \\
\midrule
\texttt{gr\_sim\_load\_scenario(s,name,params,n)} & build a registered C scenario & \S\ref{sec:pipeline} \\
\texttt{gr\_scenario\_register(def)} & register one & --- \\
\texttt{gr\_scenarios\_init()} & populate the registry & --- \\
\bottomrule
\end{tabular}}

%====================================================================
\section{The Scenario Layer and Web Front End}
\label{sec:scenario}
%====================================================================

A run's initial conditions can be assembled two ways, both built on the API of
\S\ref{sec:api}. The native path is the C \emph{scenario registry}
(\texttt{gr\_sim\_load\_scenario}): each test in the suite is a registered
scenario plus an analytic assertion, which is how the physics of Parts~I--II is
regression-checked.

The browser path is declarative: a JSON object ---
\texttt{\{program, format, version, grid, global, background, particles, view\}}
--- is the single source of truth, and the front end's \texttt{applyScenario}
routine maps it onto exactly the calls documented above:
\begin{center}
\begin{tabular}{@{}L{5.0cm}L{8.6cm}@{}}
\toprule
\textbf{JSON field} & \textbf{API calls} \\
\midrule
\texttt{grid}                & \texttt{gr\_sim\_create} (W,H,dx,cEff,cfl) \\
\texttt{global.gEff/kE/\dots}  & couplings + scheme selectors (\S\ref{sec:api}) \\
\texttt{global.switches.*}    & the physics-switch setters (one-to-one) \\
\texttt{global.absorber/init} & friction taper + the \texttt{init\_potentials\_*} calls \\
\texttt{background[\,]}        & \texttt{gr\_sim\_clear\_background} + \texttt{gr\_sim\_add\_background\_body} per body \\
\texttt{particles[\,]}        & \texttt{gr\_sim\_add\_particle} + \texttt{set\_particle\_spin/drive}, self-field \\
\texttt{view}                & front-end visualization only (\S\ref{sec:viz}) \\
\bottomrule
\end{tabular}
\end{center}
The JSON is thus a thin, validated, shareable wrapper over the same C entry
points; a scenario fully and reproducibly determines a run. A debug bridge
additionally exposes the API over a local socket, returning \emph{computed}
diagnostics (field statistics, conservation residuals, per-particle state) so the
simulation can be driven and checked programmatically.

%====================================================================
\section{Conclusion}
%====================================================================

The three parts form one chain of custody for the model. Part~I states the
physics --- linearized GR in GEM form coupled to electromagnetism, with
relativistic, spin, and proper-time dynamics --- deriving each result in 3D and
specializing it to the $2{+}1$ dimensions the code runs in. Part~II discretizes
those 2D equations into an explicit Yee-lattice PIC scheme: leapfrog field
advance, shape-function deposition, Esirkepov charge conservation,
Lewis--Birdsall force gather, a relativistic Boris push, and the initial-field,
boundary, background, and visualization machinery. Part~III ties that scheme to
the C library, member by member and function by function. The cross-references
run the full length: an equation label in Part~I, an algorithm step in Part~II,
and a function in Part~III name the same object at three levels of description.
That is the intent --- for a simulator that borders on research work, the work is
shown.

\bigskip
\noindent\textit{End of report.}

%====================================================================
\begin{thebibliography}{9}
\addcontentsline{toc}{section}{References}
%====================================================================
\bibitem{Will1993} C.~M.~Will, \textit{Theory and Experiment in Gravitational
  Physics}, Cambridge University Press (1993), \S6.2.
\bibitem{AliHaimoud2019} Y.~Ali-Ha\"imoud, lecture notes on post-Newtonian
  theory (EIH equation of motion, eqs.~37, 41), 2019.
\bibitem{CohenMashhoon1993} J.~M.~Cohen and B.~Mashhoon, ``Standard clocks,
  interferometry, and gravitomagnetism,'' \textit{Phys.\ Lett.\ A} \textbf{181},
  353 (1993).
\bibitem{Mashhoon2003} B.~Mashhoon, ``Gravitoelectromagnetism: a brief review,''
  in \textit{The Measurement of Gravitomagnetism}, arXiv:gr-qc/0311030 (2003).
\bibitem{Iorio2001} L.~Iorio, ``Satellite gravitational orbital perturbations
  and the gravitomagnetic clock effect,'' \textit{Int.\ J.\ Mod.\ Phys.\ D}
  \textbf{10}, 465 (2001).
\bibitem{Tartaglia2000} A.~Tartaglia, ``Detection of the gravitomagnetic clock
  effect,'' \textit{Class.\ Quantum Grav.} \textbf{17}, 783 (2000).
\end{thebibliography}

\end{document}
